% =============================================================
%  Module Description — Introduction to Informatics (CS101-F26)
%  Shown on https://wosm.academy as the course page for CS101-F26.
%  Compile with: pdflatex cs101-module-description.tex
%  Write "TODO(owner): ..." where information is still missing;
%  the website highlights it.
% =============================================================
\documentclass[11pt,a4paper]{article}

\usepackage[utf8]{inputenc}
\usepackage[T1]{fontenc}
\usepackage{lmodern}
\usepackage[margin=2.2cm]{geometry}
\usepackage{booktabs}
\usepackage{tabularx}
\usepackage{array}
\usepackage{enumitem}
\usepackage{xcolor}
\usepackage{titlesec}
\usepackage{fancyhdr}
\usepackage[hidelinks]{hyperref}

% ---------- Module details (edit these) ----------------------
% \ModuleCode must equal the course's `code` in src/content/courses/.
\newcommand{\ModuleTitle}{Introduction to Informatics}
\newcommand{\ModuleCode}{CS101-F26}
\newcommand{\Programme}{Foundation programme}
\newcommand{\Level}{Semester 1}
\newcommand{\Semester}{Fall 2026}
\newcommand{\Credits}{TODO(owner)}
\newcommand{\ModuleLeader}{Dr. Ali Khudiyev}
\newcommand{\Email}{TODO(owner)}
\newcommand{\OfficeHours}{TODO(owner)}
\newcommand{\Prerequisites}{None}

% ---------- Styling ------------------------------------------
\definecolor{accent}{RGB}{0,70,130}
\titleformat{\section}{\large\bfseries\color{accent}}{\thesection.}{0.5em}{}
\titlespacing*{\section}{0pt}{1.2em}{0.5em}
\setlist{itemsep=2pt, topsep=4pt}
\setlength{\parindent}{0pt}
\setlength{\parskip}{4pt}

\pagestyle{fancy}
\fancyhf{}
\lhead{\small \ModuleCode\ -- \ModuleTitle}
\rhead{\small Module Description}
\cfoot{\small \thepage}
\renewcommand{\headrulewidth}{0.4pt}

% =============================================================
\begin{document}

\begin{center}
  {\LARGE\bfseries\color{accent} \ModuleTitle}\\[4pt]
  {\large Module Description}
\end{center}
\vspace{6pt}

% ---------- Summary table (not shown on the website: the course page
% shows these details from the macros above) --------------------------
\begin{tabularx}{\textwidth}{>{\bfseries}l X >{\bfseries}l X}
\toprule
Module code   & \ModuleCode    & Credits      & \Credits \\
Programme     & \Programme     & Level        & \Level \\
Semester      & \Semester      & Prerequisites & \Prerequisites \\
Module leader & \ModuleLeader  & Email        & \href{mailto:\Email}{\Email} \\
Office hours  & \multicolumn{3}{l}{\OfficeHours} \\
\bottomrule
\end{tabularx}

% ---------- Sections (shown on the website) ------------------
\section{Module Overview}
Introduction to Informatics (ItI) is the first-semester lecture course introducing
computer science and object-oriented programming. Each week the lecture introduces
the concepts that are then supported with tests, quizzes, and other supplementary
materials. The language of instruction and implementation is Java.

\section{Aims}
\begin{itemize}
  \item TODO(owner): aims of the module.
\end{itemize}

\section{Learning Outcomes}
After completing this course, you will:
\begin{enumerate}[label=LO\arabic*.]
  \item Understand the essential concepts of computer science and be able to apply
        them in a practical, scientifically grounded, and responsible way.
  \item Learn key concepts such as algorithms, syntax and semantics, and how to
        evaluate program efficiency in terms of memory use and runtime.
  \item Solve well-defined algorithmic problems and develop basic concurrent and
        distributed applications using Java or a similar object-oriented programming
        language. This foundation will help you learn other imperative and
        object-oriented programming languages independently.
\end{enumerate}

\section{Syllabus}
\begin{tabularx}{\textwidth}{c X}
\toprule
\textbf{Unit} & \textbf{Topic} \\
\midrule
1  & Introduction \\
2  & Control Structure \\
3  & Data Types \\
4  & Object Orientation I \\
5  & Object Orientation II \\
6  & Object Orientation III \\
7  & Algorithms \\
8  & Programming Languages \\
9  & Graphical User Interfaces \\
10 & Recursion \\
11 & Syntax, Semantics, Verification \\
12 & Beyond Programming \\
13 & Course Review \& Exam \\
\bottomrule
\end{tabularx}

\section{Teaching and Learning Methods}
The module is self-paced. All materials are provided on Moodle.

\begin{tabularx}{0.6\textwidth}{X r}
\toprule
\textbf{Activity} & \textbf{Hours} \\
\midrule
Lectures            & TODO(owner) \\
Labs / Tutorials    & TODO(owner) \\
Independent study   & TODO(owner) \\
\midrule
\textbf{Total}      & \textbf{TODO(owner)} \\
\bottomrule
\end{tabularx}

\section{Assessment}
Your knowledge will be assessed throughout the course with quizzes, tutor and
homework exercises, team project, and the final exam.

TODO(owner): weight of each assessment component and the learning outcomes it covers.

\section{Reading List}
\begin{itemize}
  \item Sedgewick, R.\ and Wayne, K.\ \textit{Computer Science: An Interdisciplinary
        Approach} (2016) and \textit{Introduction to Programming in Java} (2017).
        Companion material:
        \href{https://introcs.cs.princeton.edu/java/home}{introcs.cs.princeton.edu}.
\end{itemize}

\section{Materials and Licence}
Course materials are published under the
\href{https://creativecommons.org/licenses/by-sa/4.0/}{Creative Commons Attribution-ShareAlike 4.0}
licence unless stated otherwise. See \href{https://wosm.academy/licence/}{wosm.academy/licence} for details.

\section{Policies}
\begin{itemize}
  \item \textbf{Attendance:} TODO(owner)
  \item \textbf{Late submissions:} TODO(owner)
  \item \textbf{Academic integrity:} All submitted work must be your own.
        See the \href{https://wosm.academy/code-of-conduct/}{Code of Conduct}.
\end{itemize}

\end{document}
